% Discussion, limitations and conclusion.

\section{Discussion}\label{sec:discussion}

\paragraph{What a practitioner should take from this.}
If the aim is long-run control of a risk threshold under prompt shift, use an online update with a calibration window of the right length. The update gives the guarantee (Theorem~\ref{thm:longrun}), and the best window differed between our streams. Among the online rules we tried, DtACI was the strongest on the synthetic streams and on MMLU. PromptShift-CRC is a reasonable choice when the right window is unknown: in every setting it did better than the worst window of the Adaptive CRC family, and it was in the better half of the family on MMLU and RealToxicityPrompts. Its computed prompt drift $D_t$ and effective size $n_{\mathrm{eff},t}$ are useful diagnostics in their own right, because they say how far the current prompt is from the evidence behind the threshold; we have not shown that they predict violations, and that is the most useful next experiment.

\paragraph{Why the weights add little.}
Theorem~\ref{thm:longrun} gives the update full responsibility for the long-run rate. The weights can then only improve the \emph{local} rate and the sharpness. Two things limit what they can do. First, a kernel on sentence embeddings groups prompts by topic, and the risk of a prompt depends only weakly on its topic, so a topic-similar past prompt is a weak guide to the current risk; MMLU is the clearest example, since the median risk differs across categories while the upper tail does not. Second, down-weighting costs effective size (Lemma~\ref{lem:m}), and with a threshold at the $0.9$ quantile that cost is paid in the level of the quantile. Where the upper tail does move with a feature that the embedding captures (RealToxicityPrompts) the weights help a little and a longer window helps as much.

\section{Limitations}\label{sec:limits}

\begin{itemize}[nosep]
  \item \textbf{Scale of the real-output experiments.} We used three small open models (Qwen2.5-1.5B-Instruct, GPT-2 and Qwen2.5-0.5B) on CPU, two toxicity classifiers, one prompt template, one sampled continuation per prompt, and one embedding model. Larger models, other templates and other embeddings may behave differently. We generated toxic text only to measure its toxicity, and we release risk scores and not the text. Only RealToxicityPrompts has an upper tail that moves; a second shifting task would strengthen the conclusions.
  \item \textbf{The synthetic generators and their constants.} The generators and the per-experiment constants (bandwidth, gate centre and slope) are those of the earlier version of this work and were chosen while the generators were being developed, so they are not independent of the generators. We did not tune any constant after seeing the results reported here, and the constants on the real streams follow rules fixed in advance (Section~\ref{sec:protocol}).
  \item \textbf{Theory.} Theorem~\ref{thm:cond} needs conditional independence of the risks given the features and applies at a fixed level; Corollary~\ref{cor:drift} needs a Lipschitz condition that we cannot check. The long-run result (Theorem~\ref{thm:longrun}) holds for any weights but says nothing about local behaviour. We have no guarantee that the online rule controls the local rate.
  \item \textbf{Proxy streams.} The streams built from TruthfulQA, RealToxicityPrompts, SummEval and LongHalQA fields in the earlier version contain no model outputs and some are ordered by a field linked to the outcome. We report them in Appendix~\ref{app:proxy} and draw no conclusion from them.
  \item \textbf{Uncertainty on the real streams.} Standard errors on the real streams are over orderings of one fixed set of items, so they measure ordering variability only, not variability across item sets or generations. The released risk scores are canonical: sampled continuations are not bit-reproducible across hardware.
  \item \textbf{The toxicity task.} The toxicity classifier is the ground truth, and $81\%$ of the risks lie below $0.01$, close to its floor; the $0.9$ quantile in the three lowest bands is in that range. Only abrupt block shifts have known boundaries, so post-shift violation is reported for them; smooth orderings (embedding direction, length) are reported with local violation only.
  \item \textbf{Weak localisation.} With the bandwidth set to the median pairwise distance, the kernel weights are close to uniform on all of the real streams, not only on RealToxicityPrompts, which limits what ``the kernel weights contribute least'' can mean. TruthfulQA multiple-choice accuracy of the small model is near chance, so that task is uninformative about weighting.
  \item \textbf{Baselines.} The density-ratio rule is a simplified stand-in for published methods, and we did not reimplement localised adaptive risk control \citep{zecchin2024localized}.
\end{itemize}

\section{Conclusion}\label{sec:conclusion}

We set out to show that drift-aware weighting improves conformal risk control for foundation models under prompt shift. What we found is narrower. The online update provides the long-run guarantee for any weights; the weights and the gate change local validity by less than one percentage point and sharpness by up to $1.6$ points; a matched short window and DtACI are as good or better in most of our settings; and PromptShift-CRC performs like a mid-ranked window without having to choose one. We state a conditional bound with a computable prompt-drift term, give an evaluation protocol for comparing such rules fairly, and release the code, the model outputs and the streams so that others can test the drift-aware idea on tasks where the upper tail of the risk moves with the prompt.
